% paper/sections/09e_ppe_bc.tex
%
% §9.6  境界条件と数値安定性
%
\subsection{境界条件と数値安定性}
\label{sec:ppe_bc_stability}

PPE に適用する境界条件の種類と，その実装上の注意点を以下にまとめる．
\begin{itemize}
    \item \textbf{固体壁・流入出：} $\partial p/\partial n = 0$（ノイマン条件）．CCD ブロック方程式への組み込み手順（境界行を $p'_0=0$ で置き換え，変密度項の消滅）は §~\ref{sec:ccd_neumann_bc} を参照．
    \item \textbf{自由表面：} $p = p_{\mathrm{atm}}$（ディリクレ条件）
    \item \textbf{全周ノイマンの場合：} 解の一意性を保つため，1点の圧力を固定するか平均値を0にする拘束を加える（§~\ref{sec:pinpoint}参照）．
    \item \textbf{周期境界：} 全辺を周期とする場合は FVM 行列に折り返しフェイスおよびゴーストノード拘束を組み込む（§~\ref{sec:fvm_periodic} 参照）．零空間対策は壁面 BC と同一（1点ピン条件）．
\end{itemize}

%% Neumann BC 実装検証セット（無負荷試験・製造解収束テスト・零モード確認）は
%% 付録~\ref{sec:ccd_neumann_unit_test} に移動．
%% FVM PPE 周期境界条件対応（折り返しフェイス・ゴーストノード拘束・ILU 充填）は
%% 付録~\ref{sec:fvm_periodic} に移動．

\subsubsection{高密度比における条件数と数値安定性}
\label{sec:ppe_condition_number}

気液界面を挟む密度比 $\rho_l/\rho_g \approx 10^3$（水/空気相当）は
PPE 演算子の条件数に直接影響する．
なお，条件数の悪化にとどまらず，界面型密度場（smoothed Heaviside）では
変密度 PPE 自体が $\rho_l/\rho_g \geq 10$ で\textbf{DC 反復が発散
（divergent residual）}するか，反復が形式収束する場合でも
\textbf{精度が $\Ord{1}$ 付近で plateau}する本質的問題が存在する
（§~\ref{sec:varrho_ppe_limitation}；
ここで「発散」は反復残差の単調増大，「精度 plateau」は
収束しても誤差ノルムが格子細分化に対して低減しない現象を指す）．
この本質的制約に対する解決策として，§~\ref{sec:split_ppe_framework} で\textbf{分相 PPE 解法}を定式化した．
以下の安定化戦略は，smoothed Heaviside 一括解法を低密度比で用いる場合の条件数緩和策である．
水--空気級の高密度比では一括解法を主戦略とせず，分相 PPE と界面ジャンプ条件を用いる．
可変係数演算子 $\mathcal{L}_h$ の係数 $a_{i+1/2} = 2/(\rho_i+\rho_{i+1})$ は，
液相（$\rho \approx \rho_l$）と気相（$\rho \approx \rho_g$）の格子フェイスで
$\Ord{\rho_l/\rho_g}$ の比率を持つため，行列の条件数は
\begin{equation}
  \kappa(\mathcal{L}_h) = \Ord{\frac{\rho_l}{\rho_g}\cdot\frac{1}{h^2}}
  \approx O\!\left(\frac{10^3}{h^2}\right)
  \label{eq:ppe_cond_number}
\end{equation}
と見積もられる．

\noindent\textbf{安定化戦略：}
\begin{itemize}
  \item \textbf{フェイス係数 $2/(\rho_L+\rho_R)$（式~\eqref{eq:af_exact}）：}
        FVM フラックス連続性から自然に導かれるフェイス係数 $a_{i+1/2} = 2/(\rho_L+\rho_R)$ は，
        密度の算術平均の逆数であると同時に比体積 $1/\rho$ の調和平均でもある．
        \textbf{2 値の場合の等価性}：
        $\bigl[(\rho_L+\rho_R)/2\bigr]^{-1}$（算術平均の逆数）
        $=$ $2/(\rho_L+\rho_R)$；
        一方 $\bigl[(1/\rho_L+1/\rho_R)/2\bigr]^{-1}$（$1/\rho$ の調和平均）
        $=$ $2\rho_L\rho_R/(\rho_L+\rho_R)$．
        $\rho_L \to \rho_R$ 近傍では両者が $\Ord{(\rho_L-\rho_R)^2}$ で一致し，
        一般の 2 値では前者が FVM フラックス連続性と整合する
        （§\ref{sec:bf_p4_beta}；n ≥ 3 値では両者は分離する）．
        この係数は界面フラックスに物理的根拠を持ち，
        条件数の過大評価を防ぐ（付録~\ref{app:fvm_face_coeff}）．
  \item \textbf{分相 PPE + DC $k{=}3$ + CCD 勾配（§~\ref{sec:split_ppe_framework}）：}
        高密度比では各相内の定数密度 PPE を DC で解き，
        速度補正の勾配に CCD を適用する．
        Balanced-Force 条件は圧力勾配と表面張力勾配に同じ CCD 演算子を使う範囲で維持される．
  \item \textbf{ゲージ固定点の対称性（§~\ref{sec:pinpoint}）：}
        ピン固定点を中心ノード $(N/2, N/2)$ に設定することで（コーナーピン禁止），
        高密度比条件下でもグローバル非対称性 $\Ord{\|\mathrm{rhs}\|}$ を回避する．

\end{itemize}

\subsubsection{全周 Neumann 境界における零モード対策}
\label{sec:ppe_zero_mode}

全境界が $\partial p/\partial n = 0$（ノイマン条件）の場合，
PPE の係数行列は特異（零固有値あり）であり，解 $p$ は定数の自由度を持つ．
この自由度を取り除かないと反復ソルバーが発散する．

\begin{enumerate}
  \item \textbf{1点固定（ピン条件，推奨）：}
        セル $(i_0, j_0)$ の圧力を $p_{i_0,j_0} = 0$ に固定し，その行を単位行ベクトルに置き換える．
        Matrix-free 反復では各ステップ後に強制上書きする：
        \begin{equation}
          (\delta p)^m_{i_0, j_0} \leftarrow 0
          \qquad \text{（全反復ステップで更新直後に適用）}
          \label{eq:pin_overwrite}
        \end{equation}
  \item \textbf{平均値拘束：} $\sum_{i,j} p_{i,j} = 0$ の線形拘束を追加する（精度的に優位だが実装が複雑）．
\end{enumerate}

\subsubsection{圧力固定点（ピン点）の選択原則}
\label{sec:pinpoint}

\noindent\textbf{原則1：固定点を途中で変更しない．}
ピン点を変更すると，圧力の基準値が不連続に変化し圧力勾配に突発的誤差が生じる．

\noindent\textbf{原則2：固定点は純粋気相または純粋液相の領域に置く．}
界面が固定点を通過すると，ピン条件により正しい PPE 方程式が反映されず収束が崩れる．
4倍対称な設定では中心セル $(N/2,N/2)$ が数値対称性を保持する
（§~\ref{sec:ppe_gauge_neumann} 参照）．

\noindent\textbf{原則3：圧力のゼロ点は相対量である．}
$p_{i_0,j_0} = 0$ は数学的拘束であり，速度補正に必要な勾配 $\bnabla p$ はスカラーシフトに依存しない．

\phantomsection\label{sec:fccd_bc_options}% backward-compat
